\documentclass[12pt]{article}

\usepackage{amsmath}
\usepackage{amssymb}
\usepackage{amsfonts}
\usepackage{geometry}
\usepackage{setspace}
\usepackage[T1]{fontenc}
\usepackage[utf8]{inputenc}

\geometry{margin=1in}
\setstretch{1.2}

\begin{document}

\subsection*{Eraker, Johannes, and Polson (2003): The Impact of Jumps in Volatility and Returns}

\subsubsection*{Introduction}

A well-documented finding in the literature is that models combining diffusive
stochastic volatility with jumps in returns --- such as the SVJ model of Bates
(1996) --- are still misspecified. The core problem is that volatility in these models
can only move gradually, driven by normally distributed diffusion shocks, yet
empirical data shows that conditional volatility can increase very rapidly around
periods of market stress. Eraker, Johannes, and Polson (2003) address this gap
by introducing jumps directly into the volatility process, arguing that such a
factor is rapidly moving, like jumps in returns, but also persistent, unlike them.
The paper provides the first likelihood-based empirical evidence for the presence
of volatility jumps using S\&P~500 and Nasdaq~100 index return data.

\subsubsection*{The Model}

The paper considers four nested models. Writing $Y_t = \log S_t$ for the
log-price and $V_t$ for the instantaneous variance, the most general dynamics are:

\begin{equation}
\begin{pmatrix} dY_t \\ dV_t \end{pmatrix}
= \begin{pmatrix} \mu \\ \kappa(\theta - V_{t-}) \end{pmatrix} dt
+ \sqrt{V_{t-}}
\begin{pmatrix} 1 & 0 \\ \rho\,\sigma_v & \sqrt{1-\rho^2}\cdot\sigma_v \end{pmatrix}
dW_t
+\begin{pmatrix} \xi^y\,dN_t^y \\ \xi^v\,dN_t^v \end{pmatrix}
\end{equation}

where $W_t$ is a standard Brownian motion in $\mathbb{R}^2$, $N_t^y$ and $N_t^v$
are Poisson processes with intensities $\lambda_y$ and $\lambda_v$, and $\xi^y$,
$\xi^v$ are the jump sizes in returns and volatility respectively. This nests four
models: the pure stochastic volatility model \textbf{SV} (no jumps), the
\textbf{SVJ} model (jumps in returns only, $\xi^v\sim N(\mu_y,\sigma_y^2)$,
$\lambda_v=0$), the \textbf{SVCJ} model (contemporaneous jumps in both returns
and volatility, $N_t^y=N_t^v=N_t$, with $\xi^v\sim\text{Exp}(\mu_v)$ and
$\xi^y|\xi^v\sim N(\mu_y+\rho_J\xi^v,\sigma_y^2)$), and the \textbf{SVIJ} model
(independent jumps in returns and volatility). The exponential distribution for
volatility jump sizes guarantees $V_t$ stays positive and captures the large
upward spikes in volatility observed empirically.

The intuition behind including volatility jumps is straightforward. Jumps in
returns generate large discrete moves but have no lasting impact on future
return distributions. Diffusive volatility is persistent but moves slowly.
Volatility jumps combine both features: they arrive suddenly and cause an
immediate, large increase in $V_t$, which then mean-reverts gradually back to
its long-run level. This gives the model a mechanism to generate the kind of
rapid, persistent volatility increases observed around crash episodes.

\subsubsection*{Estimation via MCMC}

The paper estimates all four models using Markov Chain Monte Carlo (MCMC),
a likelihood-based approach that simultaneously recovers parameters, latent spot
volatilities, jump times, and jump sizes from the return time series. This goes
beyond standard calibration methods, which only estimate parameters and treat
volatility as unobserved. The MCMC posterior distribution is:

\begin{equation}
p(\Theta, J, \xi^y, \xi^v, V \mid Y)
\;\propto\; p(Y \mid \Theta, J, \xi^y, \xi^v, V)\, p(\Theta, J, \xi^y, \xi^v, V)
\end{equation}

where $J$, $\xi^y$, $\xi^v$, and $V$ are vectors of jump times, jump sizes,
and spot volatilities over the sample. The algorithm iteratively draws from the
conditional posteriors of each block, producing samples that fully quantify
parameter and volatility estimation uncertainty. Model comparison is carried out
using log-Bayes factors, which penalise model complexity automatically and
require no appeal to asymptotic distributions.

\subsubsection*{Empirical Results and Option Pricing Implications}

Estimated on daily S\&P~500 returns from 1980 to 1999, the SV model is
severely misspecified: generating the $-23\%$ crash return of October 1987
requires an almost eight standard deviation diffusion shock. Adding jumps in
returns (SVJ) improves the fit but introduces clustered jump arrivals during the
crash week --- three estimated jumps in a single week --- which contradicts the
i.i.d.\ Poisson assumption. Log-Bayes factors strongly favour the models with
volatility jumps: the log-Bayes factor comparing SV to SVIJ is $59.45$ for the
S\&P and $56.03$ for the Nasdaq, constituting overwhelming evidence. In the
SVCJ model, the 1987 crash is explained by a single return jump of $-14\%$
combined with a volatility jump that pushed $V_t$ from $40\%$ to over $50\%$,
a much more economically coherent picture than the SVJ model's cluster of jumps.

For option pricing, jumps in returns steepen the slope of the implied volatility
curve, while volatility jumps additionally raise implied volatility for deep
in-the-money options, producing the characteristic hook or ``tipping at the end''
effect documented by Duffie, Pan, and Singleton (2000). The SV model generates
a nearly flat implied volatility curve because it produces negligible conditional
non-normalities at short horizons. Estimation risk also matters: spot volatility
uncertainty dominates option price uncertainty for short-dated at-the-money
options, while parameter uncertainty dominates for out-of-the-money options and
long maturities.

\subsubsection*{Strengths, Limitations, and Relevance to QQQ Options}

The paper's main contribution is providing direct empirical evidence, from
return data alone, that both types of jumps are necessary components of a
well-specified model. The MCMC approach is a genuine methodological advance
because it recovers the full time series of latent states, making it possible to
attribute specific market episodes to specific model components. The use of long
samples spanning the 1987, 1997, and 1998 stress episodes is especially valuable
for identifying rare jump events that shorter option-based studies miss.

The main limitations follow those of the simpler jump-diffusion models. Jump
intensities $\lambda_y$ and $\lambda_v$ are kept constant, so the model cannot
capture the time-varying crash risk that is visible in option markets. The constant
intensity assumption is directly contradicted by the paper's own finding that the
jump component varies strongly over time, echoing the conclusion of Carr and Wu
(2003). Calibration using only return data also means that risk premia embedded
in option prices are not directly estimated, which limits the model's immediate
use for option pricing without additional assumptions.

For our QQQ project, the Eraker et al.\ framework provides important context.
The Nasdaq~100 results in the paper show that Nasdaq volatility is systematically
higher and more volatile than S\&P volatility, with jumps arriving roughly three
times more frequently. This is consistent with the technology-heavy composition
of the QQQ and reinforces the need for a model that allows both return and
volatility jumps. The finding that jumps in volatility, not diffusive stochastic
volatility, are the primary drivers of market stress episodes suggests that option
pricing models applied to QQQ should allow for stochastic and jump-driven
volatility to avoid the systematic mispricing of short-dated out-of-the-money
puts that simpler models produce.

\end{document}